\documentclass{ximera}

\input{../../preamble.tex}

\title{Core Analysis}

\begin{document}

\begin{abstract}
Part 2
\end{abstract}
\maketitle




The Core Functions are our basic, beginning function building blocks. Linear combinations and linear compositions create the Elementary Functions.

The Core Functions cannot be broken down any further and their analysis is generally straightforward.

As a starting point for function analysis, let's just list characteristics of our Core Functions. A full analysis would also provide reasoning.



\begin{explanation} \textbf{\textcolor{red!75!green}{Exponential Functions}} 


\[ E(x) = r^x  \, \text{ where } \, r > 0  \, \text{ and } \, r \ne 1 \]



\textbf{Domain:} The natural domain is $(-\infty, \infty)$.

\textbf{Zeros:} Core exponential functions have no zeros.



\textbf{Continuity:} Core exponential functions are continuous.



\textbf{End-Behavior:} 



If $r > 1$, core exponential functions tend to $0$ as the domain becomes unbounded negatively. Core exponential functions become unbounded positively as the domain becomes unbounded positively.



\[
\lim\limits_{x \to \infty} r^x = \infty
\]


\[
\lim\limits_{x \to -\infty} r^x = 0
\]


If $0 < r < 1$, then this is reversed.



\[
\lim\limits_{x \to \infty} r^x = 0
\]


\[
\lim\limits_{x \to -\infty} r^x = \infty
\]



People will say "the function approaches infinity as the domain appraoaches infinity" or "the function tends to infinity as the domain tends to infinity".


\textbf{Behavior (Inc/Dec):} 


If $r > 1$, then core exponential functions are increasing functions.

If $0 < r < 1$, then core exponential functions are decreasing functions.





\textbf{Global Maximum and Minimum:} Core exponential functions have no global maximum or minimum.

\textbf{Local Maximum and Minimum:}  Core exponential functions have no local maximums or minimums.

\textbf{Range:} The range of a core exponential function is $(0,\infty)$.

\end{explanation}











\begin{explanation} \textbf{\textcolor{red!75!green}{Core Logarithmic Functions}} 


\[ L(x) = \log_r(x)  \]



\textbf{Domain:} The natural domain is $(0, \infty)$.

\textbf{Zeros:} Core logarithmic functions have $1$ as their only zero.



\textbf{Continuity:} Core logarithmic functions are continuous.



\textbf{End-Behavior:} Since the domain is $(0, \infty)$, core logarithmic functions only have end-behavior as the domain become unbounded positively.


If $r > 1$, then

\[
\lim\limits_{x \to \infty} \log_r(x) = \infty
\]


The other side is called \textbf{singularity behavior}.

\[
\lim\limits_{x \to 0^+} \log_r(x) = -\infty
\]





If $0 < r < 1$, then

\[
\lim\limits_{x \to \infty} \log_r(x) = -\infty
\]


The other side is called singularity behavior.

\[
\lim\limits_{x \to 0^+} \log_r(x) = \infty
\]







\textbf{Behavior (Inc/Dec):} 




If $r > 1$, then core logarithmic functions are increasing functions.


If $0 < r < 1$, then core logarithmic functions are decreasing functions.





\textbf{Global Maximum and Minimum:} Core logarithmic functions have no global maximum or minimum.

\textbf{Local Maximum and Minimum:}  Core logarithmic functions have no local maximums or minimums.

\textbf{Range:} The range of any core logarithmic function is $(-\infty,\infty)$.

\end{explanation}

















\begin{explanation} \textbf{\textcolor{red!75!green}{Sine and Cosine}} 


\[ S(x) = \sin(x)  \]

\[ C(x) = \cos(x)  \]



\textbf{Domain:} The natural domain of both functions is $(-\infty, \infty)$.

\textbf{Zeros:} Both functions have zeros and since both functions are periodic, they have an infinite number of zeros.


\begin{itemize}
\item Zeros of $\sin(x)$ are ${k \pi \, | \, k \in \mathbb{Z}}$.
\item Zeros of $\cos(x)$ are ${\frac{\pi}{2} + k \pi\, | \, k \in \mathbb{Z}}$.
\end{itemize}


\textbf{Continuity:} Both functions are continuous.



\textbf{End-Behavior:} Both functions continue to oscillate.






\textbf{Behavior (Inc/Dec):} Both functions switch back-and-forth between increasing and decreasing with a period of $2\pi$.

\begin{itemize}
\item $\sin(x)$ increases on $\left(0, \frac{\pi}{2} \right)$
\item $\sin(x)$ decreases on $\left(\frac{\pi}{2}, \frac{3\pi}{2} \right)$
\item $\sin(x)$ increases on $\left(\frac{3\pi}{2}, 2\pi \right)$
\end{itemize}

\begin{itemize}
\item $\cos(x)$ decreases on $(0, \pi)$
\item $\cos(x)$ increases on $(\pi. 2\pi)$
\end{itemize}







\textbf{Global Maximum and Minimum:} Both functions have a global maximum of $1$ and a global minimum of $-1$.




\textbf{Local Maximum and Minimum:}  Both functions have $1$ and $-1$ as their local extrema.




\textbf{Range:} The range of both functions is $[-1, 1]$.

\end{explanation}



















\begin{explanation} \textbf{\textcolor{red!75!green}{Absolute Value Function}} 


\[ Abs(x) = |x|  \]



\textbf{Domain:} The natural domain is $[0, \infty)$.

\textbf{Zeros:} $0$ is the only zero.




\textbf{Continuity:} The core absolute value function is continuous.

\textbf{End-Behavior:} The core absolute value function becomes unbounded positively on both sides.




\[
\lim\limits_{x \to -\infty} |x| = \infty
\]


\[
\lim\limits_{x \to \infty} |x| = \infty
\]






\textbf{Behavior (Inc/Dec):} 

The core absolute value function decreases on $(-\infty, 0)$ and increases on $(0, \infty)$







\textbf{Global Maximum and Minimum:} The core absolute value function has $0$ as its global minimum and it has no global maximum.




\textbf{Local Maximum and Minimum:}  The core absolute value function has $0$ as its only local minimum and it has no local maximum.



\textbf{Range:} The range is $[0, \infty)$.

\end{explanation}





The Core Functions are our basic building blocks.  We will begin building the Elementary Functions as linear combinations of the Core Functions as well as through composition with linear functions. 


The Chain Rule will be our tool for obtaining the behavior of composite functions.



\begin{theorem} \textbf{\textcolor{-red!75!green}{Chain Rule}} 

\[ inc \circ inc = inc \]

\[ inc \circ dec = dec \]

\[ dec \circ inc = dec \]

\[ dec \circ dec = inc \]

\end{theorem}

If we will break down compositions into the their core component functions, then the analysis is much smoother.























\begin{center}
\textbf{\textcolor{green!50!black}{ooooo-=-=-=-ooOoo-=-=-=-ooooo}} \\

more examples can be found by following this link\\ \link[More Examples of Visual Behavior]{https://ximera.osu.edu/csccmathematics/precalculus2/precalculus2/visualBehavior/examples/exampleList}

\end{center}





\end{document}
